% GENERATED FILE. Edit canonical lessons, then run npm run generate:books.
% canonical-source-hash: fnv1a64:4689e0cfa2f4bc83
% canonical-lessons: HI-C127-malumnahin, HI-C127-bilkulnahin, HI-C127-sachmuch, HI-C127-achhaji, HI-C127-chalega, HI-R127-second-pass-repair, HI-R127-second-pass-body-mood

\chapter{Not at All}
\label{ch:hi-not-at-all}
\hlchaptermodality{127}

\begin{chapteropening}
\textbf{By the end of this chapter:} \emph{I can answer with don't know, not at all, truly, is that so, and that'll do.}

Complete HI-R127-second-pass-repair: repair and short answers, from cold.
\end{chapteropening}

\section[mālūm nahīṁ]{\hi{मालूम} \hi{नहीं} (mālūm nahīṁ) — don't know}

\label{lesson:HI-C127-malumnahin}



\begin{quote}
Before the new one: say the Hindi for I don't know, then the Hindi for one more time.
\end{quote}

\subsection*{You'll want to know: \hi{मालूम} \hi{नहीं}}
\textbf{\hi{मालूम} \hi{नहीं}} (\emph{mālūm nahīṁ}) — \textquotedblleft{}don't know\textquotedblright{}, \textquotedblleft{}no idea\textquotedblright{}.

\textbf{\hi{मालूम}} is \textquotedblleft{}known\textquotedblright{}. It says what \textbf{\hi{मुझे} \hi{नहीं} \hi{पता}} says, with a different word for knowing.

\begin{grammarlens}[title={What you've built}]
A second I don't know.
\end{grammarlens}

\subsection*{Your turn}
\begin{itemize}
\raggedright
  \item \emph{Say it:} \emph{mālūm nahīṁ}
  \item \emph{Say it:} \emph{mālūm nahīṁ}, once more
  \item \emph{Say it:} say \emph{mālūm nahīṁ}
  \item \emph{Recall:} say the Hindi for I don't know, then the Hindi for one more time, then say \emph{mālūm nahīṁ} again
\end{itemize}

\begin{tcolorbox}[breakable,colback=teal!4,colframe=teal!35!black,title={Before you move on}]
What does \hi{मालूम} \hi{नहीं} mean? (\textquotedblleft{}Don't know\textquotedblright{}.) Say it once more.
\end{tcolorbox}
\section[bilkul nahīṁ]{\hi{बिलकुल} \hi{नहीं} (bilkul nahīṁ) — not at all}

\label{lesson:HI-C127-bilkulnahin}



\begin{quote}
Before the new one: say the Hindi for one more time, then the Hindi for don't know.
\end{quote}

\subsection*{You'll want to know: \hi{बिलकुल} \hi{नहीं}}
\textbf{\hi{बिलकुल} \hi{नहीं}} (\emph{bilkul nahīṁ}) — \textquotedblleft{}not at all\textquotedblright{}, \textquotedblleft{}absolutely not\textquotedblright{}.

\textbf{\hi{बिलकुल}} is the \textquotedblleft{}absolutely\textquotedblright{} you know; with \textbf{\hi{नहीं}} it turns a no all the way down.

\begin{grammarlens}[title={What you've built}]
The strongest no.
\end{grammarlens}

\subsection*{Your turn}
\begin{itemize}
\raggedright
  \item \emph{Say it:} \emph{bilkul nahīṁ}
  \item \emph{Say it:} \emph{bilkul nahīṁ}, once more
  \item \emph{Say it:} say \emph{bilkul nahīṁ}
  \item \emph{Recall:} say the Hindi for one more time, then the Hindi for don't know, then say \emph{bilkul nahīṁ} again
\end{itemize}

\begin{tcolorbox}[breakable,colback=teal!4,colframe=teal!35!black,title={Before you move on}]
What does \hi{बिलकुल} \hi{नहीं} mean? (\textquotedblleft{}Not at all\textquotedblright{}.) Say it once more.
\end{tcolorbox}
\section[sacmuc]{\hi{सचमुच} (sacmuc) — truly, really}

\label{lesson:HI-C127-sachmuch}



\begin{quote}
Before the new one: say the Hindi for don't know, then the Hindi for not at all.
\end{quote}

\subsection*{You'll want to know: \hi{सचमुच}}
\textbf{\hi{सचमुच}} (\emph{sacmuc}) — \textquotedblleft{}truly, really\textquotedblright{}.

It opens on the \textbf{\hi{सच}} you know. \textbf{\hi{सचमुच}?} — \textquotedblleft{}really?\textquotedblright{}, a stronger \textbf{\hi{सच} \hi{में}?}.

\begin{grammarlens}[title={What you've built}]
Truly.
\end{grammarlens}

\subsection*{Your turn}
\begin{itemize}
\raggedright
  \item \emph{Say it:} \emph{sacmuc}
  \item \emph{Say it:} \emph{sacmuc}, once more
  \item \emph{Say it:} ask \emph{sacmuc?}
  \item \emph{Recall:} say the Hindi for don't know, then the Hindi for not at all, then say \emph{sacmuc} again
\end{itemize}

\begin{tcolorbox}[breakable,colback=teal!4,colframe=teal!35!black,title={Before you move on}]
What does \hi{सचमुच} mean? (\textquotedblleft{}Truly, really\textquotedblright{}.) Say it once more.
\end{tcolorbox}
\section[acchā jī]{\hi{अच्छा} \hi{जी} (acchā jī) — is that so?; all right then}

\label{lesson:HI-C127-achhaji}



\begin{quote}
Before the new one: say the Hindi for not at all, then the Hindi for truly.
\end{quote}

\subsection*{You'll want to know: \hi{अच्छा} \hi{जी}}
\textbf{\hi{अच्छा} \hi{जी}} (\emph{acchā jī}) — \textquotedblleft{}is that so?\textquotedblright{}, or \textquotedblleft{}all right then\textquotedblright{}, by tone.

\textbf{\hi{अच्छा}} is the \textquotedblleft{}good\textquotedblright{} you know; alone it already means \textquotedblleft{}oh, I see\textquotedblright{}. \textbf{\hi{जी}} makes it polite.

\begin{grammarlens}[title={What you've built}]
A polite way to show you are listening.
\end{grammarlens}

\subsection*{Your turn}
\begin{itemize}
\raggedright
  \item \emph{Say it:} \emph{acchā jī}
  \item \emph{Say it:} \emph{acchā jī}, once more
  \item \emph{Say it:} say \emph{acchā jī}
  \item \emph{Recall:} say the Hindi for not at all, then the Hindi for truly, then say \emph{acchā jī} again
\end{itemize}

\begin{tcolorbox}[breakable,colback=teal!4,colframe=teal!35!black,title={Before you move on}]
What does \hi{अच्छा} \hi{जी} mean? (\textquotedblleft{}Is that so?; all right then\textquotedblright{}.) Say it once more.
\end{tcolorbox}
\section[calegā]{\hi{चलेगा} (calegā) — that'll do, that works}

\label{lesson:HI-C127-chalega}



\begin{quote}
Before the new one: say the Hindi for truly, then the Hindi for is that so.
\end{quote}

\subsection*{You'll want to know: \hi{चलेगा}}
\textbf{\hi{चलेगा}} (\emph{calegā}) — \textquotedblleft{}that'll do\textquotedblright{}, \textquotedblleft{}that works\textquotedblright{}, literally \textquotedblleft{}it will go\textquotedblright{}.

It is \textbf{\hi{चलना}}, \textquotedblleft{}to go\textquotedblright{}, in the future. Offered a choice, \textbf{\hi{चलेगा}} accepts it without fuss.

\begin{grammarlens}[title={What you've built}]
Don't know, not at all, truly, is that so, that'll do.
\end{grammarlens}

\subsection*{Your turn}
\begin{itemize}
\raggedright
  \item \emph{Say it:} \emph{calegā}
  \item \emph{Say it:} \emph{calegā}, once more
  \item \emph{Say it:} say \emph{calegā}
  \item \emph{Recall:} say the Hindi for truly, then the Hindi for is that so, then say \emph{calegā} again
\end{itemize}

\begin{tcolorbox}[breakable,colback=teal!4,colframe=teal!35!black,title={Before you move on}]
What does \hi{चलेगा} mean? (\textquotedblleft{}That'll do, that works\textquotedblright{}.) Say it once more.
\end{tcolorbox}
\section[(dialogue)]{When you are lost, and short answers — a second pass}

\label{lesson:HI-R127-second-pass-repair}



\begin{quote}
You did not understand. Say so, and ask for it once more.
\end{quote}

\begin{grammarlens}[title={Grammar Lens: two ways to say I don't know}]
\noindent
\begin{tabularx}{\linewidth}{@{}>{\raggedright\arraybackslash}X>{\raggedright\arraybackslash}X@{}}
\toprule
\textbf{phrase} & \textbf{literally} \\
\midrule
\textbf{\hi{मुझे} \hi{नहीं} \hi{पता}} & to me, no knowing \\
\textbf{\hi{मालूम} \hi{नहीं}} & not known \\
\bottomrule
\end{tabularx}

Both keep the knower out of the subject slot, the way Hindi does for wanting.
\end{grammarlens}

\subsection*{Your turn}
Say these aloud:

\begin{itemize}
\raggedright
  \item ask what someone said; ask for it one more time
  \item say you don't know, two ways
  \item say \textquotedblleft{}not at all\textquotedblright{}; then \textquotedblleft{}truly?\textquotedblright{}
  \item show you are listening politely; accept a plan with \textquotedblleft{}that'll do\textquotedblright{}
\end{itemize}

\begin{tcolorbox}[breakable,colback=teal!4,colframe=teal!35!black,title={Before you move on}]
Once more? (\textbf{\hi{फिर} \hi{से}}.) Didn't understand? (\textbf{\hi{समझ} \hi{नहीं} \hi{आया}}.) What did you say? (\textbf{\hi{क्या} \hi{कहा}?}) I don't know? (\textbf{\hi{मुझे} \hi{नहीं} \hi{पता}}.) One more time? (\textbf{\hi{एक} \hi{बार} \hi{और}}.) Don't know? (\textbf{\hi{मालूम} \hi{नहीं}}.) Not at all? (\textbf{\hi{बिलकुल} \hi{नहीं}}.) Truly? (\textbf{\hi{सचमुच}}.) Is that so? (\textbf{\hi{अच्छा} \hi{जी}}.) That'll do? (\textbf{\hi{चलेगा}}.)
\end{tcolorbox}
\section[(dialogue)]{What the body needs, and how you feel — a second pass}

\label{lesson:HI-R127-second-pass-body-mood}



\begin{quote}
You have walked all afternoon. Say you are hungry, and thirsty.
\end{quote}

\begin{grammarlens}[title={Grammar Lens: things that come to you}]
\noindent
\begin{tabularx}{\linewidth}{@{}>{\raggedright\arraybackslash}X>{\raggedright\arraybackslash}X@{}}
\toprule
\textbf{phrase} & \textbf{literally} \\
\midrule
\textbf{\hi{नींद} \hi{आ} \hi{रही} \hi{है}} & sleep is coming \\
\textbf{\hi{पसीना} \hi{आ} \hi{रहा} \hi{है}} & sweat is coming \\
\textbf{\hi{मज़ा} \hi{आया}} & fun came \\
\bottomrule
\end{tabularx}

Sleep, sweat and fun all \emph{come} to you, on \textbf{\hi{आना}}.
\end{grammarlens}

\subsection*{Your turn}
Say these aloud:

\begin{itemize}
\raggedright
  \item say you are sleepy; say you have a headache
  \item say you are sweating a lot
  \item ask a nervous friend to stay calm; say you feel restless
  \item say you don't feel like it today; then that you enjoyed it
  \item a friend offers a plan: \textquotedblleft{}is that so?\textquotedblright{} then \textquotedblleft{}that'll do\textquotedblright{}
\end{itemize}

\begin{tcolorbox}[breakable,colback=teal!4,colframe=teal!35!black,title={Before you move on}]
Hungry? (\textbf{\hi{भूखा}}.) Thirsty? (\textbf{\hi{प्यासा}}.) Sleepy? (\textbf{\hi{नींद} \hi{आ} \hi{रही} \hi{है}}.) Sweat? (\textbf{\hi{पसीना}}.) A headache? (\textbf{\hi{सिरदर्द}}.) Calm? (\textbf{\hi{शांत}}.) Nervousness? (\textbf{\hi{घबराहट}}.) Restless? (\textbf{\hi{बेचैन}}.) Not in the mood? (\textbf{\hi{मन} \hi{नहीं} \hi{है}}.) Fun? (\textbf{\hi{मज़ा}}.)
\end{tcolorbox}
